\chapter{Conclusões}
\label{cap:conclusoes}

Esta dissertação partiu do Trabalho de Graduação de Wayne de Paula, que
determinou os estados de polarização de ondas gravitacionais na teoria
bimétrica de Visser com o formalismo de Newman--Penrose. Vinte anos
depois, os limites sobre a massa do gráviton diminuíram mais de uma
ordem de grandeza, as PTAs detectaram um fundo de ondas gravitacionais e
a teoria da gravitação massiva ficou mais bem compreendida. O trabalho
revisou os resultados do TG com um tratamento exato e os levou até o
observável de uma PTA. As respostas às três questões da
Introdução são resumidas a seguir.

\section{Respostas às questões de pesquisa}

\paragraph{Q1. O que permanece correto no TG?} A contagem de seis
polarizações do modelo linear de Visser está correta e vale exatamente,
para qualquer razão $f/f_g$: as seis componentes da matriz de marés são
independentes. O sexto modo é o traço da perturbação. Em um termo de
massa que não é de Fierz--Pauli, ele é um escalar fantasma, de mesma
massa que os demais modos no caso de Visser. Esse é o preço de o modelo
evitar a descontinuidade vDVZ na ordem linear. Em Fierz--Pauli, a teoria
linear consistente, restam cinco polarizações. Os modos \emph{breathing}
e longitudinal passam a ser manifestações de uma única amplitude, com
$p_l=-(f_g/f)^2p_b$. O desvio dessa relação no modelo de Visser é
proporcional ao traço, o que fornece um critério para distinguir as duas
teorias.

As expressões NP do TG estão corretas como contrações do Riemann. Sua
interpretação como amplitudes de polarização e a classificação E(2),
porém, pressupõem ondas quase nulas, com erro de ordem
$\frac12(f_g/f)^2$. No caso de baixa frequência do artigo de 2004, esse
erro é de cerca de 60\%. Para ondas massivas, a separação em helicidades
depende do observador, porque o grupo que preserva o vetor de onda é
SO(3), e não E(2). Há também uma ambiguidade de fator dois na
normalização da massa no TG, que deve ser fixada pela relação de
dispersão.

A conclusão do TG sobre as bandas de frequência sobrevive em forma mais
precisa. Para componentes métricas comparáveis, as marés não tensoriais
são suprimidas por $(f_g/f)^2$ em relação às tensoriais, lei que os
números do artigo de 2004 seguem. Essa normalização, contudo, não é
fixada pela cinemática: com a normalização canônica, a parte transversal
da helicidade zero não é suprimida. O que é robusto é que as assinaturas
\emph{especificamente massivas} são suprimidas por potências de $f_g/f$.
Essas assinaturas são a dispersão, a componente longitudinal e o vínculo
entre as deformações escalares. Com os limites atuais, elas só são de
ordem um na banda das PTAs.

\paragraph{Q2. Como os modos aparecem numa PTA?} O desvio de frequência
dos pulsos, incluindo o movimento dos observadores, depende da onda
apenas pela matriz de marés. É a projeção de $E_{ij}$ sobre a linha de
visada, na Terra e no pulsar, dividida pelo fator de Doppler. O objeto
calculado no TG é, portanto, exatamente o que determina a resposta de
uma PTA. A correlação tensorial passa de Hellings--Downs a um quadrupolo
puro quando a frequência se aproxima do corte. A correlação da
helicidade zero de Fierz--Pauli passa da forma \emph{breathing}
$(3+\cos\zeta)/24$ ao mesmo quadrupolo. Ela contém termos cruzados entre
as respostas transversal e longitudinal, cuja omissão pode inverter o
sinal da correlação. No limiar, cada estado de polarização produz a
mesma correlação angular, consequência da simetria de rotação do
referencial de repouso da onda. Perto do limiar, os termos dos pulsares
não se cancelam na média angular, e as correlações passam a depender
das distâncias.

\paragraph{Q3. Quanto uma PTA pequena aprende?} Muito pouco, nas
configurações simuladas. Com 12 a 16 pulsares, 4,5 anos de dados e ruído
conhecido, a informação média sobre a massa foi de 0,07 a 0,11 nat, o
equivalente a excluir cerca de 10\% do intervalo a priori. Comprimir as
correlações em frequência reduz essa informação a cerca de 5\%. Nesse
regime, aproximações da resposta, como a avaliação em uma única
frequência de referência, parecem inofensivas apenas porque não há
informação para distorcer. Limites superiores próximos de $h/T$ podem
refletir sobretudo a priori. A aproximação normal usual para
estimadores de correlação produz subcobertura sistemática. Numa busca
de helicidade zero, ela multiplicou por cerca de dez a taxa de falsas
detecções, enquanto o poder real de detecção era baixo.

\section{Contribuições}

\begin{enumerate}
 \item Revisão exata, em $f/f_g$, da cinemática de polarização do TG,
 com a identificação do sexto modo como o traço (fantasma escalar) e a
 identidade $p_l+(f_g/f)^2p_b=-\frac{\Delta}{4}H$ que o relaciona ao
 vínculo de Fierz--Pauli.
 \item Delimitação quantitativa da validade das amplitudes NP e das
 classes E(2) para ondas massivas, e correções de normalização e sinal
 nas expressões do TG.
 \item Reinterpretação da comparação entre bandas de frequência,
 separando o que é cinemático do que depende da emissão (vDVZ e
 Vainshtein).
 \item Expressão da resposta de PTA diretamente em termos da matriz de
 marés; correlações dos setores tensorial e de helicidade zero de
 Fierz--Pauli com termos dos pulsares; explicação da degenerescência no
 limiar por simetria; identificação de uma discrepância numa expressão
 publicada.
 \item Evidência numérica, em redes pequenas, de que a compressão em
 frequência elimina quase toda a informação sobre a massa e de que a
 aproximação normal dos estimadores aumenta as falsas detecções de
 polarização escalar.
\end{enumerate}

\section{Limitações}

A análise é linear e em torno de Minkowski. A emissão por fontes
astrofísicas e o mecanismo de Vainshtein não foram tratados; por isso,
a potência da helicidade zero é um parâmetro livre. Os modos vetoriais
de Fierz--Pauli não foram incluídos nas correlações. As simulações usam
redes pequenas, ruído conhecido, modelo periódico sem ajuste de
temporização e distâncias conhecidas. Seus resultados são estudos de
método e não restrições observacionais. O regime pouco informativo das
campanhas limita, em particular, as conclusões sobre aproximações da
resposta.

\section{Trabalhos futuros}

\begin{enumerate}
 \item \textbf{Previsão realista de sensibilidade.} Estimar, para uma
 rede com dezenas de pulsares e 15--20 anos de dados, a menor massa e a
 menor fração de helicidade zero distinguíveis da RG, com ruído e
 distâncias marginalizados. É o passo necessário para saber se as
 assinaturas dos Capítulos~\ref{cap:polarizacoes} e~\ref{cap:pta} são
 mensuráveis.
 \item \textbf{Setor vetorial.} Completar as correlações das cinco
 polarizações de Fierz--Pauli e verificar a degenerescência no limiar
 prevista pelo argumento de simetria.
 \item \textbf{Amplitudes informadas pela emissão.} Substituir o
 parâmetro livre $\varepsilon$ por previsões de emissão em teorias com
 Vainshtein, ligando a potência escalar à massa.
 \item \textbf{Dados reais.} Incorporar o ajuste de temporização, a
 amostragem irregular e a incerteza nas distâncias. Aplicar o modelo aos
 dados públicos de PTAs, com verossimilhança testada no regime
 relevante.
 \item \textbf{Literatura.} Comunicar aos autores a discrepância
 identificada na expressão da correlação escalar.
\end{enumerate}
